problem:  
Let \( X \) be a smooth Fano manifold with \( \omega_0 \in c_1(X) \) and consider the normalized Kähler–Ricci flow  
\[
\frac{\partial \omega_t}{\partial t} = -\operatorname{Ric}(\omega_t) + \omega_t, \qquad \omega_{t=0} = \omega_0.
\]  
For each \( t \), denote \( \omega_t = \omega_0 + i\partial\bar{\partial}\varphi_t \) with normalization \(\int_X e^{-\varphi_t} \omega_0^n = \int_X \omega_0^n\).  
Assume that \( \{\varphi_t\} \) remains uniformly bounded in \(C^0\).  

It is known that the Kähler–Ricci flow may subconverge (after diffeomorphisms) to a weak Kähler–Einstein metric on a \(\mathbb{Q}\)-Fano variety. On the other hand, the notion of an **optimal destabilizing test configuration** \( (\mathcal{X}, \mathcal{L}) \) of \( X \)—one minimizing the normalized Donaldson–Futaki invariant over all filtrations—has been developed in the works of Blum–Jonsson and others.  

**Question:**  
If \(X\) is not K-polystable, does uniform \(C^0\)-boundedness of the Kähler–Ricci potential force the flow to degenerate (in the Gromov–Hausdorff or Cheeger–Gromov sense) to the central fiber of the *optimal destabilizing test configuration* of \(X\)?  
Equivalently, is the analytic limit of the Kähler–Ricci flow canonically determined by the unique minimizing filtration realizing the least Donaldson–Futaki weight?

Why is it a "good" problem:  
This problem refines a central open question at the interface between complex differential geometry and algebraic geometry: whether analytic limits of the Kähler–Ricci flow intrinsically encode the algebro-geometric instability of Fano manifolds. It asks for a precise correspondence between boundedness conditions in the parabolic Monge–Ampère flow and the algebraic minimizer of stability invariants (Donaldson–Futaki functionals). In doing so, it aims to uncover whether geometric flows naturally “select” the canonical degeneration predicted by K-stability theory.  

The formulation is simple—bounded potential along a normalized flow—but its resolution would require unifying analytic convergence phenomena (Kähler–Ricci flow, weak metric limits, Gromov–Hausdorff topology) with deep algebro-geometric structures (optimal destabilizing filtrations, Donaldson–Futaki minimization). This question probes the heart of the Yau–Tian–Donaldson framework and could lead to a canonical analytic construction of moduli compactifications for Fano manifolds, thereby bridging nonlinear PDE theory and algebraic geometry in a fundamental way.